% N02 | MINUTA DE RODAPÉ — Duas fronteiras possíveis; um marco histórico e um sinal correlacionável.
% Na bibliografia de apoio estão o texto de 2000, a proposta do grupo e a decisão
% institucional. Ao desenvolver a escala latouriana, distinga o alcance de uma rede, o
% vestígio que deixa e a distribuição de poder no coletivo.
% O conteúdo abaixo é editável e entra na compilação.
Em \enquote{The Anthropocene}, Crutzen e Stoermer sugeriram demarcar o início do antropoceno no fim do século XVIII, aproximando a máquina a vapor de James Watt dos primeiros sinais de aumento de certos gases de efeito estufa nos testemunhos de gelo. Em 2023, o Anthropocene Working Group propôs a base em 1952: o aumento de plutônio proveniente dos testes nucleares, registrado nos sedimentos do lago Crawford, permitiria correlacionar camadas do tempo da Grande Aceleração. Em 2024, a IUGS e a ICS confirmaram a rejeição da época formal; decidir se um vestígio serve para datar camadas é uma questão diferente de decidir quem responde pela história que o produziu.
